% CODEX REPORT::FRONTIER-JOINT1::2026-10-07
\documentclass[11pt]{article}
\usepackage[margin=1in]{geometry}
\usepackage{amsmath,amssymb,booktabs,hyperref}
\title{FRONTIER-JOINT1: Charged SAT merges across clause boundaries}
\author{WRAPPER implementation experiment}
\date{7 October 2026}
\begin{document}
\maketitle
\paragraph{Status and target.}
\textbf{COMPUTATIONAL; general compiler OPEN.} This run implements an explicit
SAT-preserving merge rule and accounts for its discovery, table construction,
source access and storage together. It does not establish a P versus NP
separation. The symbolic invariant below is supplied for review; independent
red-team promotion is pending. The construction is ordinary frontier dynamic
programming, not a claim of a new general SAT algorithm.

\section{Exact construction and invariant}
Let $F=C_1\land\cdots\land C_m$, with variables labelled $1,\ldots,n$,
$L$ literal occurrences, encoded bit length $N$, and clause width at most $q$.
Fix the input clause order; discovering a better order is not free.
Define the boundary after clause $i$ by
\[
 B_i=\operatorname{vars}(C_1,\ldots,C_i)
       \cap\operatorname{vars}(C_{i+1},\ldots,C_m),\qquad
 w=\max_i|B_i|.
\]
The table $D_i$ has one Boolean entry per assignment to $B_i$:
\[
 D_i(a)=1\iff \exists b\text{ on prefix variables}:
 b|_{B_i}=a\ \land\ \bigwedge_{j\le i}C_j(b)=1.
\]
Initially $B_0$ is empty and $D_0(\emptyset)=1$.
For a transition let $U=B_{i-1}\cup\operatorname{vars}(C_i)$.
Enumerate assignments $x$ to $U$; if $D_{i-1}(x|_{B_{i-1}})=1$
and $C_i(x)=1$, set $D_i(x|_{B_i})=1$.
The terminal table is on the empty boundary.

\paragraph{Why the merge is sound.}
Every prefix variable outside $B_i$ is absent from the suffix. Thus two
satisfying prefix histories that agree on $B_i$ admit exactly the same suffix
extensions. Their provenance can be discarded for the SAT decision objective.
The transition is exact by induction: an accepted old table entry supplies a
prefix witness; new variables in the current clause have not appeared in the
prefix unless already on $B_{i-1}$, so the witnesses glue consistently.
Conversely, every satisfying prefix witness restricts to an enumerated $x$
and a previously accepted entry. At $i=m$ the invariant is exactly SAT($F$).
This is stronger than merging histories merely because each has answer SAT.
No parent-scale SAT query or oracle is used to acquire the boundary.
It does not preserve an additional external provenance requirement unless
that requirement is also included in the state.

\section{Three charged implementations}
\paragraph{Full rescanning.}
At every cut, scan the whole source once for each variable label to decide
whether it occurs on both sides. This costs at most $mnL$ source-literal reads
plus $L$ clause reads. No occurrence index is retained.

\paragraph{Occurrence index.}
One source scan constructs first and last occurrence arrays with $2(n+1)$
integer slots. Each cut tests all labels against those arrays. Source reads are
at most $2L$ and index probes at most $2mn$. The index consumes
$O(n\log(m+2))$ bits. This is a real cost, even if read-only source input is
excluded from workspace.

\paragraph{Local rescanning.}
Only variables in $U$ can belong to the next boundary. Scan the suffix for
each variable in $U$ and retain exactly those that occur again. This avoids
both the full occurrence index and scans for unrelated variables.
For a completed run, source reads obey
\[
 R\le L+\sum_{i=1}^m |U_i|L_{>i}
 \le L+m(w+q)L.
\]
The implementation intentionally scans each suffix fully; early stopping is
not credited. A cap of 12 boundary variables is tested before allocating
the next table. Cap exhaustion returns UNKNOWN, never UNSAT.

\section{Joint resource ledger}
All versions use the same transition rule. For completed runs,
\[
 E=\sum_i2^{|U_i|}\le m2^{w+q},\qquad
 \text{simultaneous table payload}\le 2^{w+1}\ \text{bytes}.
\]
The table payload is not total Python memory. Lists of variable labels,
positions, clause signs, counters and transient assignments must also be
charged. In a conventional compact representation, the local-rescan design
uses
\[
 S=O\bigl(2^w+(w+q)\log(N+2)\bigr)\quad\text{bits},
\]
excluding the read-only input. The constant eight bits per table byte is
absorbed in this expression. This representation bound is not a Python RSS
measurement.

On an indexed word-RAM with $w+q=O(\log N)$, a conservative operation bound is
\[
 T=O\bigl(mL(w+q)+L+m2^{w+q}(w+q)^2\bigr).
\]
The squared factor safely charges projection and linear searches in the small
variable lists. General larger $w$ incurs multiword integer overhead; the
formula is not asserted as a unit-cost bound for unbounded integers.
A conversion to a particular lower-bound machine model must be supplied
before claiming a time-space contradiction.

For fixed $q$ and $w=o(\log N)$, the local-rescan representation has
$N^{o(1)}$ work bits and polynomial runtime. This is an achievable restricted
region, not a uniform guarantee for arbitrary CNF. In particular, the source
scan term can be quadratic in input size up to logarithmic factors.
For the chain below, the first variable remains live until the final unit
clause; scanning the successively shorter suffixes already makes
$\Omega(n^2)$ literal reads, although $w=2$. Since $N=\Theta(n\log n)$ for
this encoding, the implemented discovery procedure cannot attain a
subquadratic target on this family. This is a lower bound on this scan
procedure only; the family itself is easy.

A boundary-state description uses $w$ assignment bits, but the algorithm
retains up to $2^w$ viability entries. The distinction prevents confusing a
short state name with a small complete table. The informal $2^k/2^x$ intuition
corresponds to forgetting variables that no future clause can distinguish;
it does not bound discovery cost or prove that all $2^k$ histories are feasible.

\section{Reproducible tests}
The batch tests 290 formulas: all 256 subsets of the eight full-width signed
clauses on three variables; 18 seeded random formulas on 6, 8 and 10 variables;
12 structured formulas on 8, 16, 32 and 64 variables; and four boundary/UNSAT
controls. Each of three variants runs on every formula.
There are 861 completed answer comparisons and nine cap exits. Each cap exit
is answered correctly by the previously verified source-scanning propagation
backend. The fallback is called after the capped solver returns, so its tables
are no longer retained. Eighteen renaming/complementation checks also pass.
No mismatches were found. The complete batch ran in about 0.73 seconds under
a 25-second alarm; no long-running chunk was needed. The JSON records the
seed, exact code SHA256, input dimensions, table payload, source reads,
index probes, transitions and fallback metrics.

\begin{center}
\begin{tabular}{lrrr}
\toprule
64-variable chain & Full scan & Local scan & Indexed\\
\midrule
Source literal reads & 520319 & 11909 & 254\\
Index probes & 0 & 0 & 8192\\
Transition assignments & 502 & 502 & 502\\
Peak table bytes & 8 & 8 & 8\\
Index integer slots & 0 & 0 & 130\\
\bottomrule
\end{tabular}
\end{center}
The chain is $\bigwedge_{v<n}(\neg x_v\lor x_{v+1})\land x_1$.
The listed values are counters, not timing ratios or total memory figures.

\paragraph{Order and easy-family control.}
For 64 independent variables, interleaving $(x_v)$ and
$(\neg x_v\lor x_v)$ yields boundary width one and three peak table bytes.
Listing all units first and all tautologies second yields width 64 and trips
the cap at the thirteenth unit. The formulas have exactly the same clause
multiset and are trivially satisfiable. The fallback uses 64 propagations
and 12480 literal reads on the delayed order. This is a hostile control for
raw syntactic width, not evidence of SAT hardness. Tautology deletion would
repair this particular example cheaply, but does not prove a generic
small-width transformation. Random instances reach width ten already at ten
variables, so the test supplies no general width-reduction guarantee.
The fallback retains a linear-length assignment trail and has no favorable
worst-case branching theorem; completeness of the combined procedure does
not preserve the restricted small-space bound.

\section{Literature and corpus reconciliation}
Allender et al. study the time-space tradeoff for width-parameterized SAT
\cite{allender2014}. Their work supports treating table space and recomputation
time together; it does not give this project a general small-width compiler.
No lower bound from their conjectural consequences is imported unconditionally.

\paragraph{Corpus reviewed and duplication audit.}
This experiment used the current root MAIN executive state, CODEX guide,
repository research skill and experiment-mode protocol, the local preflight
snapshot, the preceding source-backed code/results, repository-wide pathwidth
search results, and the full GROK H1719 and UNIVERSALDAG1 records.
The preflight succeeded on a partial local snapshot, not a full checkout.
This was a targeted experiment review, not a fresh exhaustive audit of all
historical source files. H1719 already documents clause-order sensitivity,
separator-exactness, and the failure of ordered width to imply unrestricted
hardness. UNIVERSALDAG1 already identifies the need to compare residual DAGs
with OBDD/FBDD and semantic quotients. These ideas are adopted, not claimed
as discoveries of this run. The new contribution is an executable charged
boundary implementation, a cheaper local rescan, and a reproducible joint
ledger with explicit cap exits.

\paragraph{What changed relative to root MAIN.}
A concrete sufficient SAT merge relation is instantiated on declared ordered
CNF inputs. No general source-certified contraction, ordinary-CNF return
compiler, arbitrary-algorithm lower bound or P versus NP separation is added.
No root MAIN claim is strengthened or retired. Main4 and main5 receive the
experiment note; executive synthesis is handed off under the ownership rule.

\paragraph{Barrier audit and transfers.}
The boundary projection is an exact dynamic-programming interface. Its
resemblance to information passing across communication cuts is a useful
analogy, not an unrestricted lower-bound transfer. The existing easy-family
controls already defeat that transfer. No cryptographic or random-oracle
assumption is introduced. Oracle access could bypass this solver's table;
there is no nonrelativizing hinge here. Algebraic-oracle resistance and
Natural-Proofs avoidance are not claimed. No circuit lower bound is inferred.

\paragraph{Next concrete target.}
Replace quadratic suffix rereading by a uniformly constructed small-space
boundary-access schedule with a charged total near-linear bound, and test
whether it applies beyond already narrow ordered formulas. Any favorable
all-input claim must also bound $w$, or replace explicit $2^w$ tables by a
sound cheaper representation. Both remain OPEN.

\bibliographystyle{plain}
\bibliography{frontier_refs}
% END_OF_VERIFIED_REPORT::CODEX::FRONTIER-JOINT1
\end{document}
